\documentclass[11pt,letterpaper]{article}
\usepackage{cornell-notes}
\usepackage{mathtools}

\title{Chapter 10: The Importance of (Sub)sequence Comparison in Molecular Biology}
\author{}
\date{}

\setDocTitle{Chapter 10: The Importance of (Sub)sequence Comparison in Molecular Biology}
\setDocSubtitle{String Algorithms}
\setDocVersion{v1.0}
\setDocStatus{Study Companion}
\setDocOwner{String Algorithms Cornell Notes}
\setDocAuthor{}
\setDocDate{\today}
\setDocSummary{Cornell-format study and review notes for Chapter 10: The Importance of (Sub)sequence Comparison in Molecular Biology.}

\setCornellCollection{String Algorithms}
\setCornellUnitType{Chapter}
\setCornellUnitNumber{10}
\setCornellUnitTitle{The Importance of (Sub)sequence Comparison in Molecular Biology}
\setCornellCompanionLabel{Study and review companion}

\hypersetup{pdftitle={Chapter 10: The Importance of (Sub)sequence Comparison in Molecular Biology},pdfsubject={Cornell notes},pdfauthor={}}

\begin{document}
\maketitle

\begin{CornellOverviewBox}{Chapter overview}
\textbf{Purpose.} Frame sequence comparison as a computational model for biological relatedness, function, structure, and mutation while keeping algorithmic evidence distinct from biological interpretation.

\textbf{Learning objectives}
\begin{itemize}
  \item Explain why approximate comparison is necessary for molecular sequences.
  \item Distinguish substring, subsequence, similarity, and homology claims.
  \item Connect mutation models to scoring choices.
  \item Recognize the inferential limits of a high-scoring alignment.
\end{itemize}

\textbf{Prerequisites.} Basic DNA, RNA, protein, exact matching, and elementary probability.
\end{CornellOverviewBox}

\begin{CornellNotesTable}
\CornellNoteRow{\S10: Why compare molecular sequences?}{Related sequences accumulate substitutions, insertions, deletions, duplications, and rearrangements. Exact equality is therefore too brittle. A comparison algorithm searches for conserved order or locally similar regions that may indicate shared ancestry, structure, or function.}
\CornellNoteRow{What is a sequence model?}{DNA and proteins are represented as strings over finite alphabets, but symbols have unequal biological meanings. A model specifies which changes are allowed and assigns costs or scores reflecting their plausibility. The algorithm optimizes within that model; it does not validate the model itself.}
\CornellNoteRow{Substring or subsequence?}{A substring is contiguous. A subsequence preserves order while permitting omitted symbols. Alignments make omissions explicit as gaps, allowing the same representation to describe insertions and deletions. The correct abstraction depends on whether contiguity is essential.}
\CornellNoteRow{Similarity or homology?}{Similarity is a measured property of an alignment or scoring scheme. Homology is a historical relationship and is not a percentage. Strong statistically significant similarity may support a homology hypothesis, but database composition, repeats, and model choice affect the evidence.}
\CornellNoteRow{Global or local evidence?}{Global alignment asks whether two complete sequences correspond end to end. Local alignment searches for high-scoring regions embedded in otherwise unrelated sequence. Domain sharing, exon structure, and recombination make local evidence especially important.}
\CornellNoteRow{What must a responsible result report?}{State the scoring matrix, gap model, boundary policy, optimum score, alignment coordinates, and statistical calibration when available. Separate algorithmic optimality under the scoring function from biological confidence.}
\CornellNoteRow{How does this motivate later algorithms?}{Edit distance models minimum changes; similarity DP maximizes evidence; affine gaps represent runs of indels; local alignment finds conserved segments; database search introduces fast filters and significance estimates; multiple alignment and trees address families rather than pairs.}
\CornellNoteRow{What should practice emphasize?}{Given a biological question, choose global versus local scope, specify operations and weights, and identify confounders. Explain what conclusion the resulting score can and cannot justify.}
\end{CornellNotesTable}

\begin{CornellSummaryBox}{Chapter synthesis}
Molecular sequence comparison is an inference pipeline: encode sequences, choose a change or score model, optimize an alignment, and interpret the result with statistical and biological context. Approximate methods are necessary because meaningful relatedness persists despite mutations, but algorithmic similarity must not be confused with proof of common ancestry or function.
\end{CornellSummaryBox}

\begin{CornellSummaryBox}{Key takeaways}
\begin{CornellRecallList}
  \item Mutation makes exact matching insufficient for most evolutionary questions.
  \item A scoring scheme is a model, not a neutral measurement.
  \item Substrings require contiguity; subsequences require order only.
  \item Global and local alignment answer different questions.
  \item Similarity is measurable; homology is a historical hypothesis.
  \item Interpretation requires coordinates, model details, and significance context.
\end{CornellRecallList}
\end{CornellSummaryBox}

\begin{CornellOverviewBox}{Algorithm selection: when to use this}
\begin{CornellChecklist}
  \item Use global comparison for sequences expected to correspond across their full lengths.
  \item Use local comparison for shared domains or embedded conserved regions.
  \item Use exact matching for known motifs when variation is not allowed.
  \item Use a calibrated database-search method for large collections and significance estimates.
\end{CornellChecklist}
\end{CornellOverviewBox}

\begin{CornellExamBox}{Self-test questions}
\begin{CornellRecallList}
  \item Why can two related sequences fail exact matching?
  \item What is the difference between a substring and a subsequence?
  \item Why is a high alignment score model dependent?
  \item When is local alignment preferable to global alignment?
  \item What additional evidence is needed before interpreting similarity biologically?
\end{CornellRecallList}
\end{CornellExamBox}

{\footnotesize
\textbf{Terms and notation to remember.} substitution, insertion, deletion, indel, substring, subsequence, alignment, similarity, homology, global alignment, local alignment.

\textbf{Connections.} The next chapters formalize edit and similarity objectives with dynamic programming, then refine them for space, speed, multiple sequences, databases, and evolutionary trees.
}

\end{document}
